% from ega0-18-fr.tex
\oldpage[0\textsubscript{IV-1}]{55}

\begin{env}[18.2.3]
\label{0.18.2.3-fr}
On dit qu'une $A$-extension $E$ de $B$ par un $B$-bimodule $L$ est
\emph{$A$-triviale} si $E$ est un $A$-anneau augmenté (sur $B$) trivial
(18.1.4). On définit sur le $A$-bimodule produit $B\times L$ une structure de
$A$-anneau en posant $(x,s)(y,t)=(xy,xt+sy)$, comme on le vérifie aussitôt,
et il est immédiat que les applications canoniques $j:L\to B\times L$ et
$f:B\times L\to B$ définissent une $A$-extension de $B$ par $L$ qui est
$A$-triviale, l'application canonique $g:B\to B\times L$ étant un
$A$-homomorphisme inverse à droite de $f$. On dit que cette extension est
\emph{l'extension triviale type} de $B$ par $L$ et on la note $D_B(L)$~; il
est immédiat que toute $A$-extension $A$-triviale de $B$ par $L$ est
$A$-équivalente à $D_B(L)$.

On notera que toute extension du $A$-anneau $A$ lui-même par un $A$-bimodule
est nécessairement $A$-triviale.
\end{env}

\begin{env}[18.2.4]
\label{0.18.2.4-fr}
Étant données deux $A$-extensions
$L\xrightarrow{j}E\xrightarrow{f}B$,
$L'\xrightarrow{j'}E'\xrightarrow{f'}B'$, un \emph{morphisme} de la première
dans la seconde est par définition un triplet d'homomorphismes de
$A$-bimodules $(u,v,w)$ tel que le diagramme
\begin{equation}
\label{0.18.2.4.1-fr}
\tag{18.2.4.1}
  \xymatrix{
    0\ar[r] & L\ar[r]^j\ar[d]_w & E\ar[r]^f\ar[d]_u
      & B\ar[r]\ar[d]_v & 0\\
    0\ar[r] & L'\ar[r]_{j'} & E'\ar[r]_{f'} & B'\ar[r] & 0
  }
\end{equation}
soit commutatif, $u$ et $v$ étant des $A$-homomorphismes \emph{d'anneaux} et
$w$ étant tel que $w(bz)=v(b)w(z)$ et $w(zb)=w(z)v(b)$ pour $z\in L$ et
$b\in B$ (en d'autres termes, le couple $(v,w)$ constitue un
\emph{di-homomorphisme} du $B$-bimodule $L$ dans le $B'$-bimodule $L'$)~; il
est clair que si $(u',v',w')$ est un morphisme de
$L'\to E'\to B'$ dans une $A$-extension $L''\to E''\to B''$,
$(u'\circ u,v'\circ v,w'\circ w)$ est encore un morphisme, ce qui justifie la
terminologie.

La considération des deux carrés commutatifs du diagramme (18.2.4.1) va nous
conduire à deux opérations sur les extensions de $A$-anneaux.
\end{env}

\begin{env}[18.2.5]
\label{0.18.2.5-fr}
En premier lieu, considérons une $A$-extension $E'$ de $B'$ par $L'$
\[
  0\longrightarrow L'\xrightarrow{j'}E'\xrightarrow{f'}B'\longrightarrow0
\]
et un $A$-homomorphisme d'anneaux $v:B\to B'$, et soit
$F=E'\times_{B'}B$ \emph{l'image réciproque} par $v$ du $A$-anneau augmenté
$E'$ (18.1.5), de sorte que l'on a un diagramme commutatif
\[
  \xymatrix{
    0\ar[r] & L_0\ar[r]\ar[d]^i\ar@{}[d]|{\wr} & F\ar[r]^{p_2}\ar[d]^{p_1}
      & B\ar[r]\ar[d]^v & 0\\
    0\ar[r] & L'\ar[r]_{j'} & E'\ar[r]_{f'} & B'\ar[r] & 0
  }
\]
dont les lignes sont exactes, $p_1$ et $p_2$ étant les homomorphismes
canoniques~; on a vu (18.1.3) que $i$ est bijectif et il résulte aussi de la
définition (18.1.2) que $L_0^2=0$, de sorte que l'on peut considérer que $F$
est une $A$-extension de $B$ par $L'$, que l'on appelle \emph{image
réciproque par $v$} de l'extension $E'$ de $B'$ par $L'$ ($L'$ étant
naturellement considéré comme $B$-bimodule au moyen de l'homomorphisme
d'anneaux $v$). Le caractère fonctoriel
